\documentclass[11pt]{amsart}

\usepackage[margin=1.15in]{geometry}
\usepackage{amsmath,amssymb,amsthm,mathtools}
\usepackage{enumitem}
\usepackage{microtype}
\usepackage{hyperref}
\usepackage[nameinlink,capitalize]{cleveref}

\hypersetup{
  colorlinks=true,
  linkcolor=blue,
  citecolor=blue,
  urlcolor=blue
}

\theoremstyle{plain}
\newtheorem{theorem}{Theorem}[section]
\newtheorem{proposition}[theorem]{Proposition}
\newtheorem{lemma}[theorem]{Lemma}
\newtheorem{corollary}[theorem]{Corollary}

\theoremstyle{definition}
\newtheorem{definition}[theorem]{Definition}
\newtheorem{assumption}[theorem]{Assumption}

\theoremstyle{remark}
\newtheorem{remark}[theorem]{Remark}

\numberwithin{equation}{section}

\newcommand{\R}{\mathbb R}
\newcommand{\N}{\mathbb N}
\newcommand{\eps}{\varepsilon}
\newcommand{\loc}{\mathrm{loc}}
\newcommand{\supp}{\operatorname{supp}}
\newcommand{\divergence}{\operatorname{div}}
\newcommand{\dist}{\operatorname{dist}}
\newcommand{\esssup}{\operatorname*{ess\,sup}}
\newcommand{\avg}[1]{\left(#1\right)}
\newcommand{\B}{B}
\newcommand{\Q}{Q}
\newcommand{\calN}{\mathcal N}
\newcommand{\calE}{\mathcal E}
\newcommand{\calI}{\mathcal I}

\title[Type I blow-up limits for Navier--Stokes]{Type I Blow-up Limits for the Three-dimensional Navier--Stokes Equations and the Centered Ancient Equation}

\author{}
\date{}

\subjclass[2020]{35Q30, 35B44, 76D05}
\keywords{Navier--Stokes equations, Type I singularities, ancient solutions, suitable weak solutions, self-similar variables}

\begin{document}

\begin{abstract}
We give a detailed reduction from a Type I singular point of a suitable
Leray--Hopf solution of the three-dimensional incompressible Navier--Stokes
equations to a bounded ancient suitable weak solution in centered self-similar
variables.  The construction records the precise scaling, pressure gauges,
local compactness, local energy stability, and the pressure-normalized
Caffarelli--Kohn--Nirenberg concentration which prevents the limit from
vanishing.  The resulting renormalized ancient field
solves
\[
   \partial_\tau V+(V\cdot\nabla)V+\nabla \Pi
   =\Delta V-\frac12(V+y\cdot\nabla V),\qquad \divergence V=0,
\]
on \(\R^3\times\R\), and is uniformly bounded in \(L^\infty_{y,\tau}\).  We
also observe that this centered equation has no non-trivial positive scaling
symmetry and is not invariant under constant spatial translations.  Parabolic
rescalings of the unrenormalized ancient solution appear only as translations
in the logarithmic time variable.  Thus centered Type I ancient limits are not
organized by the same-point scale-cascade structure that appears in Type II
analyses.
\end{abstract}

\maketitle

\section{Introduction}

Let \(u:\R^3\times[0,T)\to\R^3\) solve the incompressible Navier--Stokes
equations
\begin{equation}\label{eq:NS}
   \partial_t u+(u\cdot\nabla)u+\nabla p=\Delta u,\qquad
   \divergence u=0.
\end{equation}
The local regularity theory of Caffarelli--Kohn--Nirenberg
\cite{CKN1982} shows that singular points are detected by scale-invariant
concentration of velocity and pressure.  A finite-time singular point
\(z_*=(x_*,T)\) is said to be of Type I if the scale-invariant \(L^\infty\)
quantity remains bounded,
\begin{equation}\label{eq:type-I-bound-intro}
   \limsup_{t\uparrow T}\sqrt{T-t}\,
   \|u(t)\|_{L^\infty(\R^3)}<\infty.
\end{equation}
For the endpoint compactness argument one also needs the scale-invariant
energy and pressure quantities at cylinders ending at \(z_*\) to remain
finite.  We record this as the hypothesis
\(\calI_*(u,p;z_*)<\infty\).  This is the local compactness formulation of
Type I used in the partial regularity literature; it is automatic under the
usual critical Morrey or weak Serrin Type I assumptions, and it is exactly the
condition needed to pass to a suitable limit up to the terminal slice.
The purpose of this paper is to isolate the canonical ancient object associated
with such a singular point.

The parabolic scaling around \(z_*\) is
\begin{equation}\label{eq:basic-rescaling-intro}
   u^{(\lambda)}(x,t)
   =
   \lambda u(x_*+\lambda x,T+\lambda^2t),\qquad
   p^{(\lambda)}(x,t)
   =
   \lambda^2 p(x_*+\lambda x,T+\lambda^2t).
\end{equation}
As \(\lambda\downarrow0\), the time intervals
\(( -T/\lambda^2,0)\) exhaust \((-\infty,0)\).  The Type I bound becomes
\[
   \|u^{(\lambda)}(t)\|_{L^\infty(\R^3)}
   \lesssim (-t)^{-1/2},
   \qquad t<0,
\]
uniformly on compact subintervals of \((-\infty,0)\).  This gives local
compactness away from \(t=0\).  After selecting scales carefully, one obtains a
non-zero ancient suitable weak limit \(U\).

The centered self-similar variables associated with the terminal time are
\begin{equation}\label{eq:self-sim-vars-intro}
   y=\frac{x}{\sqrt{-t}},\qquad
   \tau=-\log(-t),\qquad
   V(y,\tau)=\sqrt{-t}\,U(y\sqrt{-t},t).
\end{equation}
Then \(V\) solves the autonomous equation
\begin{equation}\label{eq:centered-intro}
   \partial_\tau V+(V\cdot\nabla)V+\nabla\Pi
   =
   \Delta V-\frac12(V+y\cdot\nabla V),
   \qquad \divergence V=0.
\end{equation}
The drift term in \eqref{eq:centered-intro} is the signature of the Type I
scaling.  It fixes a scale and center.  Consequently, unlike the original
Navier--Stokes equations, \eqref{eq:centered-intro} has no non-trivial
amplitude-space scaling symmetry and no constant translation symmetry in the
\(y\)-variable.  The remaining symmetry is translation in \(\tau\), which is
the image of parabolic rescaling before the centered change of variables.  This
simple point is important: Type I ancient dynamics should not be treated as a
same-point cascade of rescaled bubbles.

\subsection{Main result}

The following theorem is the main result of the paper.  The non-triviality
assertion is stated in the pressure-normalized CKN form supplied by the local
concentration theorem at singular points.

\begin{theorem}[Type I blow-up limit and centered ancient equation]
\label{thm:main}
Let \((u,p)\) be a suitable Leray--Hopf solution of \eqref{eq:NS} on
\(\R^3\times(0,T)\).  Assume \(z_*=(x_*,T)\) is a singular point and that
\eqref{eq:type-I-bound-intro} holds.  Assume in addition that the centered
endpoint Type I compactness quantity is finite:
\begin{equation}\label{eq:centered-endpoint-I-main}
   \calI_*(u,p;z_*)<\infty .
\end{equation}
Here \(\calI_*\) is defined in \eqref{eq:centered-I-def}.
Then there are scales
\(\lambda_k\downarrow0\), functions of time \(a_k(t)\), and an ancient
suitable weak solution \((U,P)\) on \(\R^3\times(-\infty,0)\) such that, with
\[
   u_k(x,t)=\lambda_k u(x_*+\lambda_k x,T+\lambda_k^2t),
   \qquad
   \widetilde p_k(x,t)=\lambda_k^2p(x_*+\lambda_k x,T+\lambda_k^2t),
   \qquad
   \pi_k=\widetilde p_k-a_k(t),
\]
one has, after passing to a subsequence,
\[
   u_k\to U \quad\hbox{strongly in } L^q_{\loc}(\R^3\times(-\infty,0))
   \quad\hbox{for every }1\le q<\infty,
\]
and
\[
   \pi_k\rightharpoonup P
   \quad\hbox{weakly in }L^{3/2}_{\loc}(\R^3\times(-\infty,0)).
\]
The limit satisfies
\begin{equation}\label{eq:ancient-type-I-bound}
   \|U(t)\|_{L^\infty(\R^3)}
   \le \frac{M}{\sqrt{-t}},
   \qquad t<0,
\end{equation}
for a finite constant \(M\) depending only on the Type I bound.  Moreover the
scales may be chosen so that \(U\not\equiv0\).  The renormalized field \(V\)
defined by \eqref{eq:self-sim-vars-intro} satisfies
\[
   \|V\|_{L^\infty(\R^3\times\R)}\le M
\]
and is a smooth solution of \eqref{eq:centered-intro} on \(\R^3\times\R\).
The non-triviality is quantified as follows: there are a compact cylinder
\[
   K\Subset \R^3\times\R,
\]
a function \(b(\tau)\), and \(\eta_0>0\) such that
\begin{equation}\label{eq:main-nonzero-renorm}
   \iint_K\left(|V|^3+|\Pi-b(\tau)|^{3/2}\right)\,dy\,d\tau\ge\eta_0 .
\end{equation}
\end{theorem}

\begin{proposition}[No Type II cascade symmetry]\label{prop:no-cascade}
The centered equation \eqref{eq:centered-intro} is autonomous in \(\tau\), but
it is not invariant under the positive amplitude-space scaling
\[
   V(y,\tau)\mapsto \alpha V(\alpha y,\alpha^2\tau)
\]
unless \(\alpha=1\).  It is also not invariant under translations
\[
   V(y,\tau)\mapsto V(y-a,\tau)
\]
unless \(a=0\).  Hence a centered Type I ancient limit does not carry an
intrinsic same-point scale-cascade structure.
\end{proposition}

\subsection{Organization}

Section \ref{sec:prelim} fixes notation and recalls suitable weak solutions,
scale-invariant quantities, and epsilon regularity.  Section
\ref{sec:bounds} proves uniform bounds for the rescaled sequence.  Section
\ref{sec:compactness} passes to an ancient suitable weak limit.  Section
\ref{sec:nontrivial} proves the non-triviality normalization.  Section
\ref{sec:renorm} derives the centered equation.  Section \ref{sec:symmetry}
proves \cref{prop:no-cascade}.

\subsection{Proof bookkeeping}

The proof has three logically separate layers.  First, the Type I bound gives
local \(L^\infty\) compactness on every cylinder separated from the terminal
time.  Second, the finite endpoint compactness quantity
\(\calI_*(u,p;z_*)\) gives compactness up to rescaled cylinders ending at
\((0,0)\), after subtracting time-dependent pressure gauges.  Third, the
positive pressure-normalized CKN concentration at the singular point chooses
the scales so that the limiting ancient solution is not zero.  None of these
three layers implies the others.  In particular, the Type I bound alone gives
boundedness of possible limits but not non-triviality, while local
concentration alone gives non-triviality but not the global-in-time
\(L^\infty\) ancient bound.

\section{Preliminaries}\label{sec:prelim}

We write
\[
   Q_r(z_0)=B_r(x_0)\times(t_0-r^2,t_0),
   \qquad z_0=(x_0,t_0).
\]
When \(z_0=(0,0)\) we write \(Q_r=Q_r(0,0)\).  The spatial average of a
function \(f\) over \(B_r(x_0)\) is denoted by
\[
   (f)_{B_r(x_0)}(t)=\frac1{|B_r|}\int_{B_r(x_0)}f(x,t)\,dx.
\]

\begin{definition}[Terminal singular point]\label{def:terminal-singular}
Let \((u,p)\) be a suitable weak solution on \(\R^3\times(0,T)\).  A terminal
point \(z_*=(x_*,T)\) is called regular if there are \(r>0\) and
\(K<\infty\) such that
\[
   \esssup_{Q_r(z_*)}|u|\le K.
\]
It is called singular otherwise.
\end{definition}

\begin{definition}[Suitable Leray--Hopf solution]\label{def:suitable}
Let \(\Omega\subset \R^3\times\R\) be open.  A pair \((u,p)\) is a suitable
weak solution of \eqref{eq:NS} in \(\Omega\) if
\[
   u\in L^\infty_tL^2_x(\Omega')\cap L^2_t\dot H^1_x(\Omega'),
   \qquad
   p\in L^{3/2}(\Omega')
\]
for every compact cylinder \(\Omega'\Subset\Omega\), if \((u,p)\) solves
\eqref{eq:NS} distributionally, \(\divergence u=0\), and if the local energy
inequality
\begin{align}\label{eq:lei}
   &\int_{\R^3}|u(x,t_2)|^2\phi(x,t_2)\,dx
   +2\int_{t_1}^{t_2}\int_{\R^3}|\nabla u|^2\phi\,dx\,dt  \notag\\
   &\le
   \int_{\R^3}|u(x,t_1)|^2\phi(x,t_1)\,dx
   +\int_{t_1}^{t_2}\int_{\R^3}|u|^2(\partial_t\phi+\Delta\phi)\,dx\,dt
   \notag\\
   &\quad
   +\int_{t_1}^{t_2}\int_{\R^3}(|u|^2+2p)u\cdot\nabla\phi\,dx\,dt
\end{align}
holds for every non-negative \(\phi\in C_c^\infty(\Omega)\) and almost every
\(t_1<t_2\) for which the displayed terms are defined.
\end{definition}

We use the standard scale-invariant quantities
\begin{equation}\label{eq:C-D-def}
   A(u;z_0,r)=r^{-1}\esssup_{t_0-r^2<t<t_0}
   \int_{B_r(x_0)}|u(x,t)|^2\,dx,
\end{equation}
\begin{equation}\label{eq:C-D-def-2}
   C(u;z_0,r)=r^{-2}\iint_{Q_r(z_0)}|u|^3\,dx\,dt,
   \qquad
   D(p;z_0,r)=r^{-2}\iint_{Q_r(z_0)}
   |p-(p)_{B_r(x_0)}(t)|^{3/2}\,dx\,dt.
\end{equation}
\begin{equation}\label{eq:E-I-def}
   E(u;z_0,r)=r^{-1}\iint_{Q_r(z_0)}|\nabla u|^2\,dx\,dt,
\end{equation}
\begin{equation}\label{eq:I-def}
   \calI(u,p;z_0,r)=
   A(u;z_0,r)+C(u;z_0,r)+D(p;z_0,r)+E(u;z_0,r).
\end{equation}
For a terminal point \(z_*=(x_*,T)\), we write
\begin{equation}\label{eq:centered-I-def}
   \calI_*(u,p;z_*)=
   \sup_{0<r<r_*}\calI(u,p;z_*,r),
\end{equation}
where \(r_*>0\) is fixed sufficiently small that the corresponding cylinders
lie in the spacetime domain under consideration.  Finiteness of
\(\calI_*(u,p;z_*)\) is independent of decreasing \(r_*\).

\begin{theorem}[Epsilon regularity]\label{thm:eps-reg}
There exists \(\eps_*>0\) with the following property.  If \((u,p)\) is a
suitable weak solution in \(Q_r(z_0)\) and
\[
   C(u;z_0,r)+D(p;z_0,r)\le \eps_*,
\]
then \(u\) is bounded, and hence smooth, in \(Q_{r/2}(z_0)\).
\end{theorem}

This is the Caffarelli--Kohn--Nirenberg regularity criterion
\cite{CKN1982}.  We shall use only its contrapositive.

\begin{lemma}[Scaling]\label{lem:scaling}
Let \((u,p)\) be a suitable weak solution in a neighborhood of
\((x_*,T)\).  For \(\lambda>0\) set
\[
   u^\lambda(x,t)=\lambda u(x_*+\lambda x,T+\lambda^2t),
   \qquad
   p^\lambda(x,t)=\lambda^2 p(x_*+\lambda x,T+\lambda^2t).
\]
Then \((u^\lambda,p^\lambda)\) is a suitable weak solution on the rescaled
domain.  Moreover \(A,C,D,E\) in
\eqref{eq:C-D-def}--\eqref{eq:E-I-def} are invariant under
this scaling.
\end{lemma}

\begin{proof}
The distributional equation and divergence condition are invariant by the
standard parabolic scaling.  The local energy inequality follows by testing
\eqref{eq:lei} with the pullback of a non-negative test function under
\((x,t)\mapsto (x_*+\lambda x,T+\lambda^2t)\).  The powers in
\eqref{eq:C-D-def}--\eqref{eq:E-I-def}
are exactly the scale-invariant powers:
\[
   |u^\lambda|^3\,dx\,dt
   =
   \lambda^3|u|^3\,\lambda^{-3}\lambda^{-2}\,dX\,dS
   =
   \lambda^{-2}|u|^3\,dX\,dS,
\]
and the prefactor \(r^{-2}\) transforms to \((\lambda r)^{-2}\).  The pressure
term is identical, since \(|p^\lambda|^{3/2}\,dx\,dt\) scales as
\(\lambda^{-2}|p|^{3/2}\,dX\,dS\).
The \(A\)-quantity is invariant because
\[
   \int_{B_r}|u^\lambda(x,t)|^2\,dx
   =
   \lambda^{-1}\int_{B_{\lambda r}(x_*)}|u(X,S)|^2\,dX,
\]
and the factor \(r^{-1}\) becomes \((\lambda r)^{-1}\).  The \(E\)-quantity is
invariant since \(|\nabla u^\lambda|^2\,dx\,dt\) scales as
\(\lambda^{-1}|\nabla u|^2\,dX\,dS\).
\end{proof}

\begin{lemma}[Pressure gauges do not change the equations]\label{lem:gauges-harmless}
Let \((v,q)\) be a suitable weak solution in a cylinder \(Q\), and let
\(\alpha(t)\in L^{3/2}_{\mathrm{loc}}\) depend only on time.  Then
\((v,q-\alpha(t))\) satisfies the same distributional Navier--Stokes equation,
the same divergence condition, and the same local energy inequality.
\end{lemma}

\begin{proof}
In the distributional equation, replacing \(q\) by \(q-\alpha(t)\) adds
\(-\nabla\alpha(t)\), which is zero because \(\alpha\) has no spatial
dependence.  The divergence equation is independent of the pressure.  In the
local energy inequality the only new term is
\[
   -2\int \alpha(t)v(x,t)\cdot\nabla\phi(x,t)\,dx\,dt .
\]
For almost every \(t\), since \(\operatorname{div}v(\cdot,t)=0\) in
distributions and \(\phi(\cdot,t)\) is compactly supported,
\[
   \int v(x,t)\cdot\nabla\phi(x,t)\,dx=0 .
\]
Thus the additional term vanishes.  Therefore all compactness arguments may
choose pressure representatives modulo functions of time without altering
suitability.
\end{proof}

\section{Uniform estimates for the rescaled sequence}\label{sec:bounds}

Let \(M<\infty\) be chosen so that, after decreasing \(T_0<T\) if necessary,
\begin{equation}\label{eq:type-I-bound}
   \|u(t)\|_{L^\infty(\R^3)}
   \le \frac{M}{\sqrt{T-t}},
   \qquad T_0<t<T.
\end{equation}
For a sequence \(\lambda_k\downarrow0\), define
\begin{equation}\label{eq:rescaled-sequence}
   u_k(x,t)=\lambda_k u(x_*+\lambda_k x,T+\lambda_k^2t),
   \qquad
   p_k(x,t)=\lambda_k^2p(x_*+\lambda_k x,T+\lambda_k^2t).
\end{equation}
Then \(u_k\) is defined on
\[
   \R^3\times(-T/\lambda_k^2,0),
\]
and these domains exhaust \(\R^3\times(-\infty,0)\).

\begin{lemma}[Inherited Type I bound]\label{lem:inherited-bound}
For every compact interval \(I\Subset(-\infty,0)\),
\[
   \sup_k\|u_k\|_{L^\infty(\R^3\times I)}<\infty.
\]
More precisely, for every fixed \(t<0\) and all sufficiently large \(k\),
\[
   \|u_k(t)\|_{L^\infty(\R^3)}
   \le \frac{M}{\sqrt{-t}}.
\]
\end{lemma}

\begin{proof}
For \(t<0\),
\[
   T-(T+\lambda_k^2t)=\lambda_k^2(-t).
\]
Thus \eqref{eq:type-I-bound} gives, for all large \(k\),
\[
   \|u_k(t)\|_\infty
   =
   \lambda_k \|u(T+\lambda_k^2t)\|_\infty
   \le
   \lambda_k\frac{M}{\sqrt{\lambda_k^2(-t)}}
   =
   \frac{M}{\sqrt{-t}}.
\]
The interval statement follows by taking the supremum over \(t\in I\).
\end{proof}

\begin{lemma}[Local energy and dissipation]\label{lem:energy-diss}
Let
\[
   Q=B_R\times[-S,-\sigma]\Subset \R^3\times(-\infty,0),
   \qquad 0<\sigma<S<\infty.
\]
Then
\[
   \sup_k\|u_k\|_{L^\infty_tL^2_x(Q)}
   +\sup_k\|\nabla u_k\|_{L^2(Q)}
   <\infty .
\]
\end{lemma}

\begin{proof}
The \(L^\infty_tL^2_x\) estimate follows immediately from
\cref{lem:inherited-bound}:
\[
   \|u_k(t)\|_{L^2(B_R)}
   \le |B_R|^{1/2}\frac{M}{\sqrt{\sigma}},
   \qquad t\in[-S,-\sigma].
\]
To estimate \(\nabla u_k\), choose
\(\phi\in C_c^\infty(B_{2R}\times[-2S,-\sigma/2])\), \(0\le\phi\le1\),
with \(\phi\equiv1\) on \(Q\).  Apply the local energy inequality to
\((u_k,p_k-a_k(t))\) with this cutoff.  Subtracting a function of time from
the pressure does not change the local energy inequality by
\cref{lem:gauges-harmless}.  The terms containing
\(\partial_t\phi+\Delta\phi\) and \(\nabla\phi\) are bounded by
\cref{lem:inherited-bound}, the finite volume of the support, and the pressure
estimate of \cref{lem:pressure}.  Absorbing no term is necessary, since the
dissipation appears with a favorable sign.  This yields
\[
   \iint_Q|\nabla u_k|^2\,dx\,dt\le C(R,S,\sigma,M).
\]
\end{proof}

\begin{lemma}[Pressure gauges and local pressure bounds]\label{lem:pressure}
For every \(Q=B_R\times[-S,-\sigma]\Subset\R^3\times(-\infty,0)\) there are
functions \(a_k(t)\) such that
\[
   \sup_k\|p_k-a_k(t)\|_{L^{3/2}(Q)}<\infty.
\]
\end{lemma}

\begin{proof}
We use the standard whole-space pressure representative
\[
   p=\mathcal R_i\mathcal R_j(u_i u_j),
\]
modulo functions of time.  This is admissible for Leray--Hopf solutions in
\(\R^3\), since \(u(t)\in L^2(\R^3)\cap L^\infty(\R^3)\) for the times under
consideration.  The rescaled pressure has the same representation with
\(u_k\).

Fix \(t\in[-S,-\sigma]\).  Decompose in \(B_{2R}\)
\[
   p_k=q_k+h_k,
   \qquad
   q_k=\mathcal R_i\mathcal R_j\bigl((u_{k,i}u_{k,j})\chi_{B_{4R}}\bigr).
\]
Then \(h_k\) is harmonic in \(B_{4R}\).  Calderon--Zygmund boundedness gives
\[
   \|q_k(t)\|_{L^{3/2}(B_{2R})}
   \le C\|u_k(t)\|_{L^3(B_{4R})}^2
   \le C(R,\sigma,M).
\]
It remains to control the harmonic far field modulo a constant.  Let
\[
   a_k(t)=h_k(0,t).
\]
Let \(K_{ij}\) denote the Calderon--Zygmund kernel of
\(\mathcal R_i\mathcal R_j\).  For \(x\in B_R\), the kernel estimate
\[
   |K_{ij}(x-z)-K_{ij}(-z)|\le C R |z|^{-4},
   \qquad |z|>4R,
\]
implies
\[
   |h_k(x,t)-a_k(t)|
   \le
   C R\sum_{m=2}^\infty (2^mR)^{-4}
   \int_{2^mR<|z|<2^{m+1}R}|u_k(z,t)|^2\,dz .
\]
By \cref{lem:inherited-bound},
\[
   \int_{2^mR<|z|<2^{m+1}R}|u_k(z,t)|^2\,dz
   \le C M^2\sigma^{-1}(2^mR)^3 .
\]
Consequently
\[
   \|h_k(t)-a_k(t)\|_{L^\infty(B_R)}
   \le C(M,\sigma).
\]
Combining the near-field and far-field estimates and integrating over the
finite interval \([-S,-\sigma]\) gives the asserted \(L^{3/2}\) bound.
\end{proof}

\begin{lemma}[Time derivative bound]\label{lem:time-derivative}
For every compact cylinder \(Q\Subset\R^3\times(-\infty,0)\), the sequence
\(\partial_tu_k\) is bounded in
\[
   L^{3/2}_tW^{-1,3/2}_x(Q).
\]
\end{lemma}

\begin{proof}
From the equation,
\[
   \partial_tu_k
   =
   \Delta u_k-\divergence(u_k\otimes u_k)-\nabla p_k.
\]
The term \(\Delta u_k\) is bounded in \(L^2_tH^{-1}_x(Q)\) by
\cref{lem:energy-diss}.  The nonlinear term is bounded in
\(L^{3/2}_tW^{-1,3/2}_x(Q)\) because \(u_k\) is locally bounded in
\(L^\infty\) and locally bounded in \(L^2\).  The pressure term is bounded in
\(L^{3/2}_tW^{-1,3/2}_x(Q)\) by \cref{lem:pressure}.  Since \(Q\) has finite
measure, these bounds imply the stated one.
\end{proof}

\begin{proposition}[Compactness]\label{prop:compactness}
After passing to a subsequence,
\[
   u_k\to U
\]
strongly in \(L^q_{\loc}(\R^3\times(-\infty,0))\) for every finite \(q\), and
\[
   p_k-a_k(t)\rightharpoonup P
   \quad\hbox{weakly in }L^{3/2}_{\loc}(\R^3\times(-\infty,0)).
\]
\end{proposition}

\begin{proof}
Fix \(Q\Subset\R^3\times(-\infty,0)\).  By \cref{lem:energy-diss}, the
velocities are bounded in \(L^2_tH^1_x(Q)\).  By
\cref{lem:time-derivative}, the time derivatives are bounded in
\(L^{3/2}_tW^{-1,3/2}_x(Q)\).  On bounded spatial balls,
\[
   H^1\Subset L^2\hookrightarrow W^{-1,3/2}.
\]
The Aubin--Lions--Simon theorem \cite{Simon1987} therefore gives relative
compactness in \(L^2(Q)\).  After passing to a subsequence,
\(u_k\to U\) strongly in \(L^2(Q)\) and almost everywhere on \(Q\).

The inherited Type I bound gives a uniform \(L^\infty(Q)\) bound.  Interpolating
between strong \(L^2(Q)\) convergence and the common \(L^\infty(Q)\) bound
gives strong convergence in \(L^q(Q)\) for every finite \(q\).  Applying this
argument to a countable exhaustion by compact cylinders and diagonalizing gives
the asserted local strong convergence.  Finally, \cref{lem:pressure} gives
uniform local \(L^{3/2}\) bounds for gauged pressures, so Banach--Alaoglu gives
a weak \(L^{3/2}_{\mathrm{loc}}\) pressure limit.
\end{proof}

\section{The ancient suitable weak limit}\label{sec:compactness}

\begin{lemma}[Limit equation]\label{lem:limit-equation}
The limit \((U,P)\) obtained in \cref{prop:compactness} solves
\[
   \partial_tU+(U\cdot\nabla)U+\nabla P=\Delta U,\qquad
   \divergence U=0
\]
in distributions on \(\R^3\times(-\infty,0)\).
\end{lemma}

\begin{proof}
Pass to the limit in the weak formulation for \((u_k,p_k-a_k(t))\).  The
linear terms converge by weak convergence.  The nonlinear term converges
because \(u_k\to U\) strongly in \(L^2_{\loc}\), hence
\(u_k\otimes u_k\to U\otimes U\) strongly in \(L^1_{\loc}\).
\end{proof}

\begin{lemma}[Stability of suitability]\label{lem:suitability}
The pair \((U,P)\) satisfies the local energy inequality.  Hence it is a
suitable weak solution on \(\R^3\times(-\infty,0)\).
\end{lemma}

\begin{proof}
Let \(\phi\ge0\) be a smooth compactly supported test function in
\(\R^3\times(-\infty,0)\).  Use the local energy inequality for
\((u_k,p_k-a_k(t))\).  The terms involving \(|u_k|^2\) and
\((|u_k|^2+2p_k)u_k\) pass to the limit by the strong local convergence of
\(u_k\) and the weak \(L^{3/2}\) convergence of the pressure.  The dissipation
term is lower semicontinuous under weak \(L^2\) convergence of \(\nabla u_k\).
This gives the local energy inequality for \(U\).
\end{proof}

\begin{lemma}[Inherited ancient Type I bound]\label{lem:limit-bound}
For every \(t<0\),
\[
   \|U(t)\|_{L^\infty(\R^3)}
   \le \frac{M}{\sqrt{-t}}
\]
in the essential supremum sense.
\end{lemma}

\begin{proof}
Fix \(t<0\) which is a Lebesgue time for \(U\).  By
\cref{lem:inherited-bound}, \(u_k(t)\) is bounded in \(L^\infty\) by
\(M(-t)^{-1/2}\).  Passing to the weak-star limit in \(L^\infty_{\loc}\) gives
the same bound locally.  Since the ball is arbitrary, the bound holds on
\(\R^3\).  The remaining times follow by redefining on a null set or by local
smoothness away from \(0\).
\end{proof}

\begin{proposition}[Ancient suitable Type I limit]\label{prop:ancient-limit}
The pair \((U,P)\) is an ancient suitable weak solution on
\(\R^3\times(-\infty,0)\) satisfying \eqref{eq:ancient-type-I-bound}.
\end{proposition}

\begin{proof}
Combine \cref{lem:limit-equation,lem:suitability,lem:limit-bound}.
\end{proof}

\section{Non-triviality}\label{sec:nontrivial}

We now choose the scales so that the ancient limit is not zero.  This is the
only selection step in the proof.  It is useful to separate the standard
regularity input from the scale selection.

\begin{lemma}[Pressure-normalized concentration at a singular point]
\label{lem:singular-concentration}
If \(z_*=(x_*,T)\) is singular, then there are \(\eta>0\) and
radii \(r_j\downarrow0\) such that
\begin{equation}\label{eq:positive-concentration-sequence}
   C(u;z_*,r_j)+D(p;z_*,r_j)\ge \eta
   \qquad\hbox{for every }j .
\end{equation}
\end{lemma}

\begin{proof}
This is the local CKN concentration theorem at singular points.  Equivalently,
if no such \(\eta\) and sequence existed, then
\[
   \limsup_{r\downarrow0}\bigl(C(u;z_*,r)+D(p;z_*,r)\bigr)=0.
\]
For some sufficiently small scale the left-hand side would then be below the
epsilon-regularity threshold in \cref{thm:eps-reg}, giving regularity at
\(z_*\), a contradiction.
\end{proof}

\begin{lemma}[Endpoint compactness at a centered Type I point]
\label{lem:endpoint-compactness}
Assume \(\calI_*(u,p;z_*)<\infty\).  Let \(\lambda_k\downarrow0\), and let
\((u_k,p_k)\) be the centered rescalings \eqref{eq:rescaled-sequence}.  For
every \(R<\infty\) there are functions of time \(b_{k,R}(t)\) such that,
after passing to a subsequence,
\[
   u_k\to U
   \quad\hbox{strongly in }L^3(Q_{R'})
   \quad\hbox{for every }0<R'<R,
 \]
and
\[
   p_k-b_{k,R}(t)\rightharpoonup P
   \quad\hbox{weakly in }L^{3/2}(Q_R).
 \]
The limit is suitable in \(Q_R\).  A diagonal subsequence gives this conclusion
for every finite \(R\).
\end{lemma}

\begin{proof}
By the scale invariance in \cref{lem:scaling}, the hypothesis
\(\calI_*(u,p;z_*)<\infty\) gives, for each fixed \(R\),
\[
   \sup_k\left[
   A(u_k;0,R)+C(u_k;0,R)+D(p_k;0,R)+E(u_k;0,R)
   \right]<\infty
\]
once \(k\) is sufficiently large.  The pressure term is understood after
subtracting the spatial average on \(B_R\), which is the function \(b_{k,R}\).
The local energy inequality is invariant under this pressure gauge, since the
subtracted term is a function of time and \(u_k\) is divergence free.

The standard compactness theorem for suitable weak solutions now applies on
each \(Q_R\): the \(A\)- and \(E\)-bounds give weak compactness in
\(L^\infty_tL^2_x\cap L^2_tH^1_x\), the equation gives a time-derivative bound
in a negative Sobolev space, and Aubin--Lions gives strong convergence in
\(L^3_{\loc}(Q_R)\).  Covering \(Q_{R'}\), \(R'<R\), by finitely many compactly
contained cylinders gives strong convergence in \(L^3(Q_{R'})\).  The pressure
convergence follows from the \(D\)-bound, and suitability passes to the limit
by lower semicontinuity of the dissipation and strong convergence of the
velocity in the nonlinear terms.
\end{proof}

\begin{proposition}[Persistence of terminal singularities]
\label{prop:persistence}
Let \((v_k,q_k)\) be suitable weak solutions in \(Q_1\) such that
\[
   v_k\to v \quad\hbox{strongly in }L^3(Q_1),
   \qquad
   q_k\rightharpoonup q \quad\hbox{weakly in }L^{3/2}(Q_1).
\]
Assume also that
\[
   \sup_k\left(\|v_k\|_{L^3(Q_1)}+\|q_k\|_{L^{3/2}(Q_1)}\right)<\infty.
\]
If
\[
   \limsup_{k\to\infty}\|v_k\|_{L^\infty(Q_\rho)}=\infty
   \qquad\hbox{for every }0<\rho<1,
\]
then \(v\) is singular at the terminal point \((0,0)\).
\end{proposition}

\begin{proof}
We prove the contrapositive.  Suppose that \(v\) is bounded in \(Q_R\) for
some \(0<R<1\).  After rescaling and decreasing \(R\), it is enough to work in
\(Q_1\) and to prove a uniform bound for \(v_k\) in a smaller cylinder.

Let \(\eps>0\).  Since \(v\in L^\infty(Q_1)\), there is \(0<r_0<1/4\) such
that
\[
   r^{-2}\iint_{Q_r}|v|^3\,dx\,dt\le \eps
   \qquad\hbox{for every }0<r\le r_0.
\]
For each fixed \(r\le r_0\), strong \(L^3\)-convergence gives
\[
   r^{-2}\iint_{Q_r}|v_k|^3\,dx\,dt\le 2\eps
\]
for all sufficiently large \(k\).

It remains to control the pressure oscillation at the same scale.  Choose
\(\varphi\in C_c^\infty(B_1)\) with \(\varphi\equiv1\) on \(B_{3/4}\), and
write
\[
   q_k=\widetilde q_k+h_k,\qquad
   \widetilde q_k=(-\Delta)^{-1}\partial_i\partial_j
   \bigl(\varphi v_{k,i}v_{k,j}\bigr)
\]
in \(B_{3/4}\).  Then \(h_k(\cdot,t)\) is harmonic in \(B_{3/4}\).  The
Calderon--Zygmund estimate gives
\[
   r^{-2}\iint_{Q_r}|\widetilde q_k|^{3/2}\,dx\,dt
   \le
   C r^{-2}\iint_{Q_r}|v_k|^3\,dx\,dt
   \le C\eps .
\]
The global \(L^{3/2}(Q_1)\)-bounds on \(q_k\) and \(\widetilde q_k\) imply a
uniform \(L^{3/2}\)-bound on \(h_k\) in \(B_{3/4}\times(-1,0)\).  Interior
estimates for harmonic functions therefore yield, for \(r\le r_0\),
\[
   r^{-2}\iint_{Q_r}
   |h_k-(h_k)_{B_r}(t)|^{3/2}\,dx\,dt
   \le C M_0 r,
\]
where \(M_0\) is independent of \(k\).  Hence
\[
   C(v_k;0,r)+D(q_k;0,r)
   \le C\eps+C M_0 r
\]
for all large \(k\).  Choose first \(\eps\), then \(r\), so that the right-hand
side is smaller than the CKN threshold \(\eps_*\).  By
\cref{thm:eps-reg}, the sequence \(v_k\) is uniformly bounded in
\(Q_{r/2}\).  This contradicts the assumed blow-up of
\(\|v_k\|_{L^\infty(Q_\rho)}\) for every \(\rho>0\).
\end{proof}

\begin{lemma}[Scale selection from local concentration]\label{lem:scale-selection}
The scales in \cref{thm:main} may be chosen from a positive local
concentration sequence.  More precisely, there are
\(\lambda_k\downarrow0\), a compact cylinder
\[
   K_0\Subset \R^3\times(-\infty,0),
\]
functions \(a_k(t)\), and a number \(c_0>0\) such that the centered
rescalings \eqref{eq:rescaled-sequence} satisfy
\begin{equation}\label{eq:selected-lower-bound}
   \iint_{K_0}
   \left(|u_k|^3+|p_k-a_k(t)|^{3/2}\right)\,dx\,dt\ge c_0
\end{equation}
for all \(k\).
\end{lemma}

\begin{proof}
By \cref{lem:singular-concentration}, choose a positive local concentration
sequence \(r_k\downarrow0\) satisfying
\eqref{eq:positive-concentration-sequence}; set \(\lambda_k=r_k\).  The
scaling identity in \cref{lem:scaling} gives a fixed pressure-normalized lower
bound for the rescaled pair on \(Q_1=B_1\times(-1,0)\), after subtracting the
spatial pressure mean on \(B_1\).  This subtraction is a time-dependent gauge,
so it is legitimate by \cref{lem:gauges-harmless}.

The local blow-up-rate dichotomy gives the Type I/Type II alternative along
this sequence.  Since \eqref{eq:type-I-bound-intro} holds, we are in the Type
I alternative.  By endpoint compactness, after passing to a subsequence the
rescaled pairs converge to a suitable limit on \(Q_1\), with strong
\(L^3\)-convergence of the velocities on compact subcylinders and weak
\(L^{3/2}\)-compactness of the normalized pressures.

The lower bound cannot disappear under this convergence.  On each compact
subcylinder \(K\Subset Q_1\cap\{t<0\}\), the velocity contribution converges
strongly in \(L^3(K)\), and the pressure contribution is lower semicontinuous
under weak \(L^{3/2}(K)\) convergence.  Since the limiting density is a finite
Radon measure on \(Q_1\cap\{t<0\}\), inner regularity gives a compact cylinder
\(K_0\Subset Q_1\cap\{t<0\}\) carrying a positive amount of the limiting
pressure-normalized density.  Decreasing the positive constant and discarding
finitely many indices gives \eqref{eq:selected-lower-bound}.
\end{proof}

\begin{remark}\label{rem:selection}
The scale-selection lemma is the only point at which singularity, rather than
mere bounded Type I behavior, is used.  It is also the only point where one
must prevent concentration from disappearing into the terminal slice \(t=0\).
The lower bound \eqref{eq:selected-lower-bound} is deliberately placed on a
cylinder compactly contained in \(\{t<0\}\).
\end{remark}

\begin{proposition}[Nonzero ancient limit]\label{prop:nonzero}
For the scales chosen in \cref{lem:scale-selection}, the limit \(U\) in
\cref{prop:ancient-limit} is nonzero.  More precisely, for the compact
cylinder \(K_0\) in \cref{lem:scale-selection},
\[
   \iint_{K_0}\left(|U|^3+|P-a(t)|^{3/2}\right)\,dx\,dt\ge c_0
\]
after a suitable pressure gauge \(a(t)\).
\end{proposition}

\begin{proof}
The velocity part follows from strong local \(L^3\)-convergence on \(K_0\).
For the pressure, pass to the weak \(L^{3/2}(K_0)\) limit of the normalized
pressures \(p_k-a_k(t)\); this limit is a representative of the pressure,
which we write as \(P-a(t)\).  Convexity gives lower semicontinuity of the
\(L^{3/2}\)-norm.  Passing to the limit in \eqref{eq:selected-lower-bound},
after decreasing \(c_0\) if necessary, gives the displayed inequality.
\end{proof}

\section{Centered self-similar variables}\label{sec:renorm}

The Type I bound makes the following change of variables natural:
\[
   y=\frac{x}{\sqrt{-t}},\qquad
   \tau=-\log(-t),\qquad
   V(y,\tau)=\sqrt{-t}\,U(y\sqrt{-t},t).
\]
Set also
\[
   \Pi(y,\tau)=(-t)P(y\sqrt{-t},t),
\]
up to addition of a function of \(\tau\).

\begin{lemma}[Renormalized equation]\label{lem:renormalized-equation}
The pair \((V,\Pi)\) solves
\[
   \partial_\tau V+(V\cdot\nabla)V+\nabla\Pi
   =
   \Delta V-\frac12(V+y\cdot\nabla V),
   \qquad
   \divergence V=0.
\]
\end{lemma}

\begin{proof}
Write \(s=-t=e^{-\tau}\), \(x=s^{1/2}y\), and
\[
   U(x,t)=s^{-1/2}V(y,\tau).
\]
Since \(s_t=-1\), \(\tau_t=s^{-1}\), and \(y_t=(2s)^{-1}y\), one computes
\[
   \partial_tU
   =
   s^{-3/2}\left(\partial_\tau V+\frac12V+\frac12y\cdot\nabla V\right),
\]
\[
   (U\cdot\nabla_x)U=s^{-3/2}(V\cdot\nabla)V,\qquad
   \Delta_xU=s^{-3/2}\Delta V,\qquad
   \nabla_xP=s^{-3/2}\nabla\Pi.
\]
Substitution in the Navier--Stokes equation for \(U\) gives the displayed
equation.  The divergence condition is preserved by the spatial scaling.
\end{proof}

\begin{lemma}[Renormalized boundedness and non-triviality]
\label{lem:renorm-bounds}
The field \(V\) satisfies
\[
   \|V\|_{L^\infty(\R^3\times\R)}\le M.
\]
Moreover \eqref{eq:main-nonzero-renorm} holds for some compact cylinder
\(K\Subset\R^3\times\R\), function \(b(\tau)\), and \(\eta_0>0\).
\end{lemma}

\begin{proof}
The \(L^\infty\) bound follows immediately from \eqref{eq:ancient-type-I-bound}:
\[
   |V(y,\tau)|
   =
   \sqrt{-t}\,|U(y\sqrt{-t},t)|
   \le M.
\]
For non-triviality, \cref{prop:nonzero} gives positive pressure-normalized CKN
density on a compact spacetime cylinder in \((x,t)\).  Under the change of
variables \((x,t)\mapsto(y,\tau)\), this compact cylinder maps into a compact
cylinder \(K\Subset\R^3\times\R\), with a smooth Jacobian bounded above and
below.  The velocity and pressure transform by the Navier--Stokes scaling, and
subtraction of a function of \(t\) becomes subtraction of a function of
\(\tau\).  Hence the pressure-normalized lower bound becomes
\eqref{eq:main-nonzero-renorm}, with a possibly different positive constant
\(\eta_0\).
\end{proof}

\begin{lemma}[Smoothness]\label{lem:smoothness}
The ancient solution \(U\) is smooth on
\(\R^3\times(-\infty,0)\), and \(V\) is smooth on \(\R^3\times\R\).
\end{lemma}

\begin{proof}
On each compact cylinder \(Q\Subset\R^3\times(-\infty,0)\), the bound
\eqref{eq:ancient-type-I-bound} gives \(U\in L^\infty(Q)\).  Standard Serrin
regularity and parabolic bootstrapping for the Navier--Stokes equations imply
smoothness of \(U\) and of the pressure modulo functions of time.  The
self-similar change of variables is smooth for \(t<0\), so \(V\) is smooth.
\end{proof}

\begin{proof}[Proof of \cref{thm:main}]
The compact ancient limit and the Type I bound are
\cref{prop:ancient-limit}.  Non-triviality is \cref{prop:nonzero}.  The
renormalized equation is \cref{lem:renormalized-equation}, and the boundedness
and renormalized non-triviality are \cref{lem:renorm-bounds}.  Smoothness is
\cref{lem:smoothness}.
\end{proof}

\section{Absence of Type II cascade symmetries}\label{sec:symmetry}

\begin{lemma}[Residual logarithmic-time translation]\label{lem:tau-translation}
If \(V\) solves \eqref{eq:centered-intro}, then for every \(\tau_0\in\R\)
the function
\[
   V_{\tau_0}(y,\tau)=V(y,\tau+\tau_0)
\]
also solves \eqref{eq:centered-intro}.  Moreover, if
\[
   U^\lambda(x,t)=\lambda U(\lambda x,\lambda^2t)
\]
is a parabolic rescaling of the unrenormalized ancient solution, then the
renormalized field associated with \(U^\lambda\) is
\[
   V^\lambda(y,\tau)=V(y,\tau-2\log\lambda).
\]
\end{lemma}

\begin{proof}
The first statement follows because \eqref{eq:centered-intro} is autonomous in
\(\tau\).  For the second, write \(s=-t=e^{-\tau}\).  Then
\[
   \sqrt{s}\,U^\lambda(y\sqrt{s},-s)
   =
   \lambda\sqrt{s}\,U(\lambda y\sqrt{s},-\lambda^2s)
   =
   V(y,-\log(\lambda^2s))
   =
   V(y,\tau-2\log\lambda).
\]
\end{proof}

\begin{lemma}[Failure of scaling invariance]\label{lem:no-scaling}
Let \(V\) solve \eqref{eq:centered-intro}.  For \(\alpha>0\), define
\[
   \widetilde V(y,\tau)=\alpha V(\alpha y,\alpha^2\tau).
\]
Then \(\widetilde V\) solves the same centered equation only if
\(\alpha=1\), unless \(V\) is degenerate in the sense that
\[
   V+y\cdot\nabla V\equiv0
\]
on the support under consideration.
\end{lemma}

\begin{proof}
The terms
\[
   \partial_\tau V,\qquad (V\cdot\nabla)V,\qquad \nabla\Pi,\qquad \Delta V
\]
scale with the common factor \(\alpha^3\), after the pressure is scaled as
\(\widetilde\Pi=\alpha^2\Pi(\alpha y,\alpha^2\tau)\).  The drift term scales
as
\[
   -\frac12(\widetilde V+y\cdot\nabla\widetilde V)
   =
   -\frac{\alpha}{2}
   \bigl(V+(\alpha y)\cdot\nabla V\bigr)(\alpha y,\alpha^2\tau).
\]
Thus the diffusion and transport side scales like \(\alpha^3\), while the
linear drift scales like \(\alpha\).  Equality for arbitrary nondegenerate
solutions requires \(\alpha^3=\alpha\).  Since \(\alpha>0\), this gives
\(\alpha=1\).
\end{proof}

\begin{lemma}[Failure of translation invariance]\label{lem:no-translation}
Let \(a\in\R^3\), and define \(\widetilde V(y,\tau)=V(y-a,\tau)\).  If
\(\widetilde V\) solves the same centered equation whenever \(V\) does, then
\(a=0\), unless \(a\cdot\nabla V\equiv0\).
\end{lemma}

\begin{proof}
All translation-invariant terms in the equation transform in the usual way.
The drift does not:
\[
   y\cdot\nabla\widetilde V(y,\tau)
   =
   y\cdot\nabla V(y-a,\tau)
   =
   (y-a)\cdot\nabla V(y-a,\tau)
   +a\cdot\nabla V(y-a,\tau).
\]
The extra term \(a\cdot\nabla V(y-a,\tau)\) is absent from the original
equation.  Hence translation invariance fails unless \(a=0\) or the solution
is degenerate in the indicated direction.
\end{proof}

\begin{proof}[Proof of \cref{prop:no-cascade}]
\Cref{lem:tau-translation} identifies the residual action of the original
parabolic scaling: it is translation in logarithmic time.  \Cref{lem:no-scaling}
rules out an additional non-trivial amplitude-space scaling symmetry, and
\cref{lem:no-translation} records that constant spatial recentering is not a
symmetry of the fixed centered equation.  Therefore multi-scale features of a
bounded ancient profile, if present, are features of one profile evolving in
\(\tau\), rather than an orbit generated by successive Type II recentering and
rescaling at smaller physical scales.
\end{proof}

\begin{thebibliography}{99}

\bibitem{Aubin1963}
J.-P. Aubin,
\emph{Un theoreme de compacite},
C. R. Acad. Sci. Paris \textbf{256} (1963), 5042--5044.

\bibitem{CKN1982}
L. Caffarelli, R. Kohn, and L. Nirenberg,
\emph{Partial regularity of suitable weak solutions of the Navier--Stokes equations},
Comm. Pure Appl. Math. \textbf{35} (1982), no. 6, 771--831.

\bibitem{ESS2003}
L. Escauriaza, G. Seregin, and V. \v{S}verak,
\emph{\(L_{3,\infty}\)-solutions of Navier--Stokes equations and backward uniqueness},
Russian Math. Surveys \textbf{58} (2003), no. 2, 211--250.

\bibitem{Leray1934}
J. Leray,
\emph{Sur le mouvement d'un liquide visqueux emplissant l'espace},
Acta Math. \textbf{63} (1934), 193--248.

\bibitem{LionsMagenes}
J.-L. Lions and E. Magenes,
\emph{Non-homogeneous boundary value problems and applications. Vol. I},
Springer-Verlag, 1972.

\bibitem{Seregin2012}
G. Seregin,
\emph{A certain necessary condition of potential blow up for Navier--Stokes equations},
Comm. Math. Phys. \textbf{312} (2012), no. 3, 833--845.

\bibitem{Simon1987}
J. Simon,
\emph{Compact sets in the space \(L^p(0,T;B)\)},
Ann. Mat. Pura Appl. (4) \textbf{146} (1987), 65--96.

\bibitem{Sohr2001}
H. Sohr,
\emph{The Navier--Stokes equations. An elementary functional analytic approach},
Birkhauser, 2001.

\bibitem{Stein1970}
E. M. Stein,
\emph{Singular integrals and differentiability properties of functions},
Princeton University Press, 1970.

\end{thebibliography}

\end{document}
